\documentclass[../../main.tex]{subfiles}

\begin{document}

\section{有界变差函数: Jordan 定理}\label{section:ch06-有界变差}

令 $f$ 为闭有界区间 $[a,b]$ 上的实值函数. 对 $[a,b]$ 的一个分割
\begin{align*}
P=\{x_{0},x_{1},\cdots,x_{k}\},\quad a=x_{0}<x_{1}<\cdots<x_{k}=b,
\end{align*}
定义分割 $P$ 上的变差 $V(f,P)$ 为
\begin{align}
V(f,P)=\sum_{i=1}^{k}|f(x_{i})-f(x_{i-1})|.\label{align:ch06-变差}
\end{align}
定义 $f$ 在 $[a,b]$ 上的全变差 $TV(f)$ 为
\begin{align}
TV(f)=\sup\{V(f,P)\mid P\text{ 是 }[a,b]\text{ 的分割}\}.\label{align:ch06-全变差}
\end{align}
我们说 $f$ 在 $[a,b]$ 上是有界变差的, 若 $TV(f)<\infty$.

若 $P'$ 是 $P$ 的加细, 则 $V(f,P')\geqslant V(f,P)$. 这是三角不等式的一个简单结论: 若在 $P$ 的两个相邻点 $x_{i-1}$ 和 $x_{i}$ 之间插入点 $y$, 则
\begin{align*}
|f(x_{i})-f(x_{i-1})|\leqslant|f(y)-f(x_{i-1})|+|f(x_{i})-f(y)|.
\end{align*}

\begin{example}\label{example:ch06-有界变差例子}
\begin{enumerate}[(1)]
\item 若 $f$ 在闭有界区间 $[a,b]$ 上是 Lipschitz 的, Lipschitz 常数为 $c$, 则它是有界变差的, 这是因为对任何分割 $P=\{x_{0},x_{1},\cdots,x_{k}\}$,
\begin{align*}
V(f,P)=\sum_{i=1}^{k}|f(x_{i})-f(x_{i-1})|\leqslant c\sum_{i=1}^{k}(x_{i}-x_{i-1})=c\cdot(b-a).
\end{align*}
\item 单调函数是有界变差的. 事实上, 若 $f$ 是递增的, 对任何分割 $P$,
\begin{align*}
V(f,P)=\sum_{i=1}^{k}|f(x_{i})-f(x_{i-1})|=\sum_{i=1}^{k}[f(x_{i})-f(x_{i-1})]=f(b)-f(a),
\end{align*}
因此 $TV(f)=f(b)-f(a)$.
\item 考虑在 $[0,1]$ 上定义的函数 $f$,
\begin{align*}
f(x)=
\begin{cases}
x^{2}\sin(\pi/x^{2})&\text{若 }x\neq 0\\
0&\text{若 }x=0
\end{cases}
\end{align*}
我们把 $[0,1]$ 分为区间 $[\frac{1}{\sqrt{n}},1]$. 对 $[0,1]$ 的分割 $P=\{0,1/\sqrt{n},1/\sqrt{n-1},\cdots,1\}$,
\begin{align*}
V(f,P)\geqslant\sum_{k=1}^{n}\left|f\left(\frac{1}{\sqrt{k+1}}\right)-f\left(\frac{1}{\sqrt{k}}\right)\right|.
\end{align*}
由于 $\sin(\pi/x^{2})$ 在 $x=1/\sqrt{k}$ 处交替取 $\pm 1$, 该和约为 $\sum_{k=1}^{n}1/k$, 它发散. 因此 $f$ 在 $[0,1]$ 上不是有界变差的.
\end{enumerate}
\end{example}

对 $[a,b]$ 的任何两个分割 $P$ 和 $Q$, 取 $P\cup Q$ 为加细, 则 $V(f,P\cup Q)\geqslant V(f,P)$ 且 $V(f,P\cup Q)\geqslant V(f,Q)$. 因此, 若 $f$ 是有界变差的, 则对任何 $c\in(a,b)$,
\begin{align}
TV(f_{[a,b]})=TV(f_{[a,c]})+TV(f_{[c,b]}).\label{align:ch06-变差可加}
\end{align}
这就是全变差的区域可加性. 我们留作习题让读者证明 (习题 31).

\begin{lemma}\label{lemma:ch06-变差引理}
令函数 $f$ 在闭有界区间 $[a,b]$ 上是有界变差的. 则 $f$ 有以下作为 $[a,b]$ 上的两个单增函数的差的直接表示:
\begin{align}
f(x)=[f(x)+TV(f_{[a,x]})]-TV(f_{[a,x]}),\quad\text{对所有 }x\in[a,b]\label{align:ch06-直接表示}
\end{align}
\end{lemma}

\begin{proof}
根据区域可加性\eqref{align:ch06-变差可加}, 对 $a\leqslant u<v\leqslant b$,
\begin{align}
TV(f_{[a,v]})-TV(f_{[a,u]})=TV(f_{[u,v]})\geqslant 0.\label{align:ch06-变差递增}
\end{align}
因此函数 $x\mapsto TV(f_{[a,x]})$, 我们称它为 $f$ 的全变差, 是一个 $[a,b]$ 上的实值的递增函数. 此外, 对于 $a\leqslant u<v\leqslant b$, 若取 $[u,v]$ 的最粗糙的分割 $P=\{u,v\}$, 我们有
\begin{align*}
f(u)-f(v)\leqslant|f(v)-f(u)|=V(f_{[u,v]},P)\leqslant TV(f_{[u,v]})=TV(f_{[a,v]})-TV(f_{[a,u]}).
\end{align*}
因此对所有 $a\leqslant u<v\leqslant b$,
\begin{align}
f(v)+TV(f_{[a,v]})\geqslant f(u)+TV(f_{[a,u]}).\label{align:ch06-和递增}
\end{align}
我们证明了以下引理.
\end{proof}

\begin{theorem}[Jordan 定理]\label{theorem:ch06-Jordan}
函数 $f$ 在闭有界区间 $[a,b]$ 上是有界变差的当且仅当它是 $[a,b]$ 上的两个单调递增函数的差.
\end{theorem}

\begin{proof}
令函数 $f$ 在 $[a,b]$ 上是有界变差的. 前一个引理提供了 $f$ 作为两个递增函数的差的直接表示. 为证明逆命题, 令 $f=g-h$, 其中 $g$ 和 $h$ 是 $[a,b]$ 上的递增函数. 对 $[a,b]$ 的任何分割 $P=\{x_{0},\cdots,x_{k}\}$,
\begin{align*}
V(f,P)&=\sum_{i=1}^{k}|f(x_{i})-f(x_{i-1})|\\
&=\sum_{i=1}^{k}|[g(x_{i})-g(x_{i-1})]+[h(x_{i-1})-h(x_{i})]|\\
&\leqslant\sum_{i=1}^{k}|g(x_{i})-g(x_{i-1})|+\sum_{i=1}^{k}|h(x_{i-1})-h(x_{i})|\\
&=\sum_{i=1}^{k}[g(x_{i})-g(x_{i-1})]+\sum_{i=1}^{k}[h(x_{i})-h(x_{i-1})]\\
&=[g(b)-g(a)]+[h(b)-h(a)].
\end{align*}
因此, $f$ 的关于 $[a,b]$ 的分割的变差的集合以 $[g(b)-g(a)]+[h(b)-h(a)]$ 为上界, 且因此 $f$ 在 $[a,b]$ 上是有界变差的.
\end{proof}

我们称将有界变差函数 $f$ 表示为两个单调递增函数的差的表达方式为 $f$ 的 Jordan 分解.

\begin{corollary}\label{corollary:ch06-BV可微}
若函数 $f$ 在闭有界区间 $[a,b]$ 上是有界变差的, 则它在开区间 $(a,b)$ 上是几乎处处可微的且 $f'$ 在 $[a,b]$ 上是可积的.
\end{corollary}

\begin{proof}
根据 Jordan 定理, $f$ 是 $[a,b]$ 上的两个递增函数的差. 于是 Lebesgue 定理告诉我们 $f$ 是两个在 $(a,b)$ 上几乎处处可微的函数的差. 因此 $f$ 在 $(a,b)$ 上几乎处处可微. $f'$ 的可积性从推论\cref{corollary:ch06-单调导数可积}得到.
\end{proof}

\begin{exercise}\label{exercise:ch06-25}
假定 $f$ 在 $[0,1]$ 上连续. 一定存在 $[0,1]$ 的非退化闭子区间 $[a,b]$, 使得 $f$ 在 $[a,b]$ 上的限制是有界变差的吗?
\end{exercise}

\begin{exercise}\label{exercise:ch06-26}
令 $f$ 为 Dirichlet 函数, 即 $[0,1]$ 中的有理数的特征函数. $f$ 在 $[0,1]$ 上是有界变差的吗?
\end{exercise}

\begin{exercise}\label{exercise:ch06-27}
在 $[0,2\pi]$ 上定义 $f(x)=\sin x$. 找出两个递增函数 $h$ 和 $g$, 使得在 $[0,2\pi]$ 上 $f=h-g$.
\end{exercise}

\begin{exercise}\label{exercise:ch06-28}
令 $f$ 为 $[a,b]$ 上的阶梯函数. 找到它的全变差的公式.
\end{exercise}

\begin{exercise}\label{exercise:ch06-29}
(i) 定义
\begin{align*}
f(x)=
\begin{cases}
x^{2}\cos(1/x^{2})&\text{若 }x\neq 0,x\in[-1,1]\\
0&\text{若 }x=0
\end{cases}
\end{align*}
$f$ 在 $[-1,1]$ 上是有界变差的吗?

(ii) 定义
\begin{align*}
g(x)=
\begin{cases}
x^{3}\cos(1/x)&\text{若 }x\neq 0,x\in[-1,1]\\
0&\text{若 }x=0
\end{cases}
\end{align*}
$g$ 在 $[-1,1]$ 上是有界变差的吗?
\end{exercise}

\begin{exercise}\label{exercise:ch06-30}
证明两个有界变差函数的线性组合也是有界变差的. 这样两个函数的积也是有界变差的吗?
\end{exercise}

\begin{exercise}\label{exercise:ch06-31}
令 $P$ 为 $[a,b]$ 的分割, 它是分割 $P'$ 的加细. 对于 $[a,b]$ 上的实值函数 $f$, 证明 $V(f,P')\leqslant V(f,P)$.
\end{exercise}

\begin{exercise}\label{exercise:ch06-32}
假设 $f$ 在 $[a,b]$ 上是有界变差的. 证明存在 $[a,b]$ 的分割序列 $\{P_{n}\}$, 使得序列 $\{TV(f,P_{n})\}$ 是递增的且收敛到 $TV(f)$.
\end{exercise}

\begin{exercise}\label{exercise:ch06-33}
令 $\{f_{n}\}$ 为 $[a,b]$ 上逐点收敛于实值函数 $f$ 的 $[a,b]$ 上的实值函数序列. 证明
\begin{align*}
TV(f)\leqslant\liminf TV(f_{n}).
\end{align*}
\end{exercise}

\begin{exercise}\label{exercise:ch06-34}
令 $f$ 和 $g$ 在 $[a,b]$ 上是有界变差的. 证明
\begin{align*}
TV(f+g)\leqslant TV(f)+TV(g)\quad\text{且}\quad TV(\alpha f)=|\alpha|\,TV(f).
\end{align*}
\end{exercise}

\begin{exercise}\label{exercise:ch06-35}
对于正数 $\alpha$ 和 $\beta$, 定义 $[0,1]$ 上的函数 $f$ 为
\begin{align*}
f(x)=
\begin{cases}
x^{\alpha}\sin(1/x^{\beta})&\text{若 }0<x\leqslant 1\\
0&\text{若 }x=0
\end{cases}
\end{align*}
证明: 若 $\alpha>\beta$, $f'$ 在 $[0,1]$ 上是可积的, 以此来证明 $f$ 在 $[0,1]$ 上是有界变差的. 接着证明: 若 $\alpha\leqslant\beta$, 则 $f$ 在 $[0,1]$ 上不是有界变差的.
\end{exercise}

\begin{exercise}\label{exercise:ch06-36}
令 $f$ 在 $[0,1]$ 上不是有界变差的. 证明存在 $[0,1]$ 中的点 $x_{0}$, 使得 $f$ 在 $[0,1]$ 的每个包含 $x_{0}$ 的长度任意小的非退化闭子区间上不是有界变差的.
\end{exercise}

\end{document}